% Separate candidate: fixed partial resolutions and fixed formal backgrounds.
% This file does not claim a tree representation count.
\subsection{Small marked sources and actual normal component presentations}

Retain the tree and marked spectral notation of
Lemmas~\ref{D:SNC-tree-residual-category} and
\ref{D:SNC-tree-minimal-marked-windows}. All scalar gerbes,
spectral covers, their formal lifts, and neutralizing lines in
this subsection are fixed geometric data. We do not assert a
bound for the coefficients of an arbitrarily chosen such lift.
The root indices and the character multiplicities may vary.

\begin{lemma}[Proved: small sources for marked spectral bundles]
\label{D:SNC-tree-small-marked-source}
Fix a component $i$, write $m=m_i$, $s=r/m$, and let
$\mathcal R$ be its marked spectral curve, with coarse curve
$R$ and reduced marked set $P$. At $x\in P$ write $d_x$
for the adjacent root order and $\mathcal P_x$ for the root
point. After one fixed positive ordinary exceptional twist,
every spectral bundle in the windows of
\ref{D:SNC-tree-minimal-marked-windows} has a presentation
\[
 0\longrightarrow L^{-1}\longrightarrow P_0
       \longrightarrow W\longrightarrow0,
 \qquad (P_0)_0=\mathcal O_R^{s+1},\qquad
 L=\det W_0.                                      \tag{D.TP.1}
\]
The kernel has weight zero at every marked point. The source
$P_0$ has an explicit all-ring presentation using only
homogeneous fibre frames at the marked points.

More precisely, put $h=s+1$ and
\[
 T=\mathcal O_{\mathcal R}
          \left(\sum_{x\in P}(d_x-1)\mathcal P_x\right).
                                                        \tag{D.TP.2}
\]
For each marked point choose the normalized source weights
$0\leq a_{x\ell}<d_x$, $1\leq\ell\leq h$, and put
$b_{x\ell}=d_x-1-a_{x\ell}$. Define the fixed sheaf
\[
 Q_x=\bigoplus_{\ell=1}^{h}
           T/T(-b_{x\ell}\mathcal P_x),
 \qquad Q=\bigoplus_{x\in P}Q_x.                    \tag{D.TP.3}
\]
A zero summand is omitted. A frame $g_x\in\operatorname{GL}_h$
defines the quotient $T^h\to Q_x$ by applying $g_x^{-1}$
and then the displayed component reductions. The kernel
\[
                 P(g)=\ker(T^h\longrightarrow Q)        \tag{D.TP.4}
\]
is a vector bundle over every coefficient algebra, has
coarse bundle $\mathcal O^h$, and realizes the flags defined
by the chosen frames and weights. These frames use
$O(ms^2)$ coordinates and quadratic inverse equations.
The construction extends over the formal component, using
the transverse root divisors in place of the root points.
After applying the fixed inverse spectral construction,
the quotient in (D.TP.1) lifts to that formal component.
Its kernel is a rank-$m$ formal bundle with residual
spectral line $L^{-1}$.
\end{lemma}

\begin{proof}
Choose the positive twist to be $-tA$, where $A$ is the
fixed strictly antinef cycle of
\ref{D:SNC-tree-uniform-Hom-depth}. Its degree on $R$ is
$tm(-A\cdot E_i)$. The lower slope window, genus and
number of marked points are all $O(m)$. A single $t$
therefore makes the coarse bundle generated and makes
the first cohomology of every Hom layer from a bundle
with coarse bundle $\mathcal O^h$ and normalized marked
weights vanish. Indeed that Hom is between $W_0^h(-P)$
and $W_0^h$; duality and the lower slope bound suffice.
In later root-normal layers the conormal degree is
nonnegative, and the marked normalization changes a
coarse degree by at most $|P|$. Enlarge $t$ once for
this further fixed correction. This is independent
of $m,s$ and every adjacent root order.

A general $s+1$-dimensional subspace of sections of the
generated rank-$s$ coarse bundle generates it: at a
point the failure of surjectivity has codimension two,
and the parameter of points has dimension one. Its
kernel is the line $L^{-1}$. At every marked point
split the surjection on fibres, lift a basis adapted
to the target flag, and supplement it by a basis of
the kernel fibre. Give the lifted vectors their
target weights and the kernel vector weight zero.
In particular the positive-weight flag spaces are
the chosen split lifts, not the full preimages of
the target spaces. The resulting strict parabolic
quotient has kernel with weight zero and
gives (D.TP.1) under the exact root/flag equivalence
\citep[Proposition3 and Theorem4]{Borne2007D}
\citep[Theorem6.1]{BorneVistoli2012D}.
Over a finite local Artinian coefficient algebra the
closed-fibre generators lift by the same $H^1$
vanishing and the coefficient-ideal filtration.
Nakayama preserves generation. With the determinant
marked, the kernel is the indicated fixed line.

We verify (D.TP.4) without assuming the coefficient
ring reduced. At $x$ take a root coordinate $z$ with
coarse coordinate $z^{d_x}$. A local generator of $T$
has character $-(d_x-1)$. The kernel of its reduction
modulo $z^{b_{x\ell}}$ is generated by
$z^{b_{x\ell}}$, and therefore has character
$-a_{x\ell}$. Conjugating the diagonal list by $g_x$
gives the asserted homogeneous free basis of $P(g)$.
Its invariant basis vectors are $g_x$ times the
ordinary basis of $\mathcal O^h$. Equivalently $Q_x$
has no invariants: its exponents lie in
$[0,b_{x\ell}-1]\subset[0,d_x-2]$, so none has
character zero after the shift $-(d_x-1)$. Exact
invariants yield $(P(g))_0=\mathcal O^h$.
The flags of this invariant bundle are exactly the
ones given by $g_x$ and the weight list. Frames and
their inverses use $2h^2$ entries at each of the
$O(m)$ points, with equations $g_xg_x^{-1}=1$.

On the lifted formal spectral cover, each marked
point is the intersection with a transverse root
divisor. Distinct such divisors are disjoint in this
formal component. Repeat (D.TP.2)--(D.TP.4) with those
divisors. Frame entries are coefficients from the
parameter algebra and hence define the same maps
on these transverse divisors. Their equations are
Cartier equations, regular even on the nilpotent
normal component. The same homogeneous kernel
calculation proves actual local freeness over every
coefficient algebra. Tensor the fixed neutralizing
line and apply finite pushforward from the fixed
degree-$m$ spectral cover. This is the inverse
spectral reference source on the formal component.

The first paragraph gives $H^1=0$ for all conormal
Hom layers from this source to the chosen positive
lattice. Thus the specified residual surjection
lifts successively through every normal layer,
including over a finite local Artinian coefficient
base. It stays surjective by Nakayama. Its kernel
is a vector bundle of rank $m$, with the stated
residual spectral line. This proves the lemma.
\end{proof}

\begin{proposition}[Proved: fixed-resolution component presentations and exact node gluing]
\label{D:SNC-tree-fixed-normal-presentation}
Fix a cycle of $f$ stages, the discrete marked weight lists, and
the formal background data specified above. Fix an
ordinary depth $a$ and a finite list of adjacent
Frobenius and bounded-pole inverse diagrams. Put
$M=\operatorname{lcm}_i(m_i)$ and $s_i=r/m_i$.
There is a polynomial presentation, with an open
subbundle and slope restriction, for actual bundle
and map data on the Cartier cycle $a\mathcal A$.
Every point over every coefficient algebra produces
actual objects. It covers the corresponding actual
families over finite local Artinian coefficient
algebras after choosing the fibre frames.

For component $i$, at each root-normal order
$1\leq j\leq aa_i m_i$ it uses $O(fm_i s_i^2)$
variables of weight $(M/m_i)j$. Its residual variables
number $O(fm_i s_i^2)$ and may all be assigned weight
$M$. Each original node contributes $O(fr^2)$
further variables of weight $M$. All equations have
weighted degree at most $KM$, where $K$ depends on
the fixed resolution, depth, characteristic and
diagram length, not on the root orders or ranks.

The formal reference data are constants, not moving
parameters. All residual determinant lines are covered.
The proposition does not by itself count coefficient
gerbes or punctured projective representations.
\end{proposition}

\begin{proof}
We first give the component construction, including
the reason that the number of marked points is not
squared. Suppress the component subscript and work
on its fixed spectral cover. As in
Lemma~\ref{D:fixed-partial-resolution-unit-atlas},
choose an ordinary fixed formal line $B$ whose
residual degree is a fixed positive multiple of $m$.
It may be the pullback of a fixed power of $-A$;
in particular it is invariant under spectral deck
translations. If $L=\det W_0$ after the positive
twist of the preceding lemma, then $|\deg L|=O(ms)$.
Choose $b=cs$, with a single fixed $c$ sufficiently
large that $L$ and $L^{-1}B^b$ have the $H^1$
vanishing and two-generator property needed below,
including the $O(m)$ marked corrections. The lower
slope bound for $W_0$ makes the assertion for $L$
uniform as well. Neither $L$ nor a basis of its
cohomology is a parameter in the construction.

Let $A_0$ be a fixed rank-$m$ formal anchor whose
residual spectral line is ordinary $\mathcal O$,
of weight zero at all adjacent root points. Such
a line is obtained by the selected neutralization;
its formal lift exists by $H^2=0$. Put
$P_j=(A_0\otimes B^{-jb})^{\oplus2}$ for $j\geq1$.
Keep the fixed objects $T$ and $Q_x$ of
(D.TP.2)--(D.TP.3), including their normal lifts;
tensor the scalar line and use inverse spectral
pushforward in those expressions when necessary.
Only their maps, including the frame matrices,
will vary. The kernel $K_1$ of the preceding
lemma has spectral line $L^{-1}$ with root weight
zero. Two generators of $L^{-1}B^b$ lift to
$P_1\twoheadrightarrow K_1$; its kernel has
spectral line $LB^{-2b}$. Continuing by two
generators gives the alternating spectral lines
\[
 K_{2v+1}=L^{-1}B^{-2vb},\qquad
 K_{2v}=LB^{-2vb}.                                  \tag{D.TP.5a}
\]
All are ordinary lines at the marked points.
At an odd step the generating Hom line is
$L^{-1}B^b$ and at an even step it is $L$.
Positive root-normal translations add degree
and alter normalized marked degrees by only
$O(m)$. Thus the same $H^1$ vanishing lifts
every specified presentation, including across
every small Artinian coefficient extension.

We obtain a partial complex of fixed terms
\[
 P_7\longrightarrow\cdots\longrightarrow P_1
            \longrightarrow T^h\longrightarrow Q,
 \qquad\text{in degrees }-7,\ldots,0,1.             \tag{D.TP.5}
\]
The last kernel is $P(g)$. Require $d^2=0$
as actual sheaf maps. For the first map
require the factored residual map
$P_1\to P(g)$ to have spectral rank at
least one; impose that condition on the
remaining maps $P_j\to P_{j-1}$ as well.
The condition is not imposed on the
displayed map $P_1\to T^h$: its composite
with the inclusion $P(g)\to T^h$ may
vanish on a marked root fibre. These rank
conditions are open on the proper residual
root curve. Locally the last quotient is
already surjective and $P(g)$ is free.
A unit minor in the preceding map splits
one rank-one summand. The next product-zero
equation and unit minor split its complement
in the rank-two term. Induction cancels the
contractible pairs. Consequently the only
cohomology is a locally free $F$ in degree
zero and a locally free end module in degree
$-7$. The same cancellation holds after any
nilpotent normal thickening and over every
coefficient ring, since each chosen unit
minor remains a unit. This proves actual
effectivity and base change; the generating
argument proves Artinian coverage. No
cohomology group depending on a flag matrix
or an unknown determinant has been inverted.

Here is a precise finite derived-Hom model for
these complexes. Compute the fixed termwise derived
Hom by the two-affine Cech complex and Cartier
resolutions of the individual $Q_x$ summands,
\[
       [\ T(-b\mathcal P_x)\xrightarrow{z_x^b}T\ ].
                                                        \tag{D.TP.6}
\]
The transverse parameter is a non-zero-divisor
even on the normal thickening and after every
coefficient-base change, so this is an actual
resolution there. Distinct supports are disjoint:
$R\mathcal Hom(Q_x,Q_y)=0$ for $x\ne y$.
Contract these zero complexes, and contract each
remaining fixed associated-graded Cech--Cartier
complex to its cohomology, over $k$. These are
$k$-linear strong contractions and remain so after
tensoring with any coefficient algebra.
The contraction is made separately for each pair
of the column objects in (D.TP.5); it preserves
their column degrees. In particular one does not
retain all cross-point terms of a direct sum of
global resolutions as independent variables.

We check the resulting sizes explicitly. At one
root point, each Hom or Ext calculation from
(D.TP.6) has coefficient exponents in an interval
of length at most $d_x$; root invariants require
one prescribed residue class modulo $d_x$.
Thus each pair of summands contributes at most a
fixed number of invariant coefficients. The same
statement holds after one $p$-Frobenius pullback,
when the intervals have length at most $pd_x$,
and after the fixed list of Cartier twists. Hence
all point-supported termwise Hom cohomology has
size $O(ms^2)$ per root-normal order. The top-top
block has multiplicity $h^2$ and line degree
$O(m)$. A block with exactly one positive-index
term has multiplicity at most $2h$ and line
degree $O(ms)$. Blocks with two such terms have
multiplicity four and line degree $O(ms)$.
Here $b=O(s)$, $\deg B=O(m)$,
$\deg_{\rm par}T=O(m)$, and the genus is $O(m)$.
There are only a fixed number of these blocks.
Curve Riemann--Roch gives
the same $O(ms^2)$ total. These are bounds for
the fixed termwise derived-Hom cohomology, not
for the much larger uncontracted Cech spaces.

Now insert the variable transitions and differentials.
The correction to the fixed differential is the
sum of positive-normal-order reference corrections
and the maps between columns of (D.TP.5). The
former raise normal order; the latter raise column
degree. A Hom between these fixed-width complexes
has column degrees between $-8$ and $8$. Since
the fixed contractions preserve column degree,
every nonzero homological-perturbation word
contains only a fixed number of variable column differentials;
all its other factors raise normal order.
Consequently the finite geometric series of
Lemma~\ref{D:weighted-fixed-residual-bundle-deformations}
terminates and supplies exact finite complexes
over every coefficient algebra. The same special
deformation-retract identities apply
\citep[Sections2.1--2.4 and3.2, complete proof]
{C_Crainic2004Perturbation}.

For clarity, the equation that $P_1\to T^h$ has
zero composite in $Q$ is an actual equality of
maps: express its class by a degree-zero cycle
and require it to be a boundary in the finite
derived-Hom complex. Since $P_1$ is locally free,
the degree-zero derived Hom to $Q$ is the
ordinary Hom, so this means exactly zero.
Keep its null homotopy in the total complex:
its lift goes from $P_1$ to the $Q$ column, of internal degree
$-1$ and column increase two. Adding this
column-skipping component makes the total
differential square to zero. This can also
be seen directly in (D.TP.6): the composite
into $T$ is divisible by its regular Cartier
equation, and the division is unique. Its
composition with $P_2$ is zero because the
Cartier equation is a non-zero-divisor and
$d_1d_2=0$. Thus there is no additional
unretained higher compatibility. Equation
(D.TP.5) denotes
this explicit twisted complex in the fixed
Cech--Cartier model when the selected chain
representatives do not compose strictly to
zero. It is quasi-isomorphic to the actual
sheaf complex whose composite is zero.
The added witness still strictly increases
column degree and obeys the same fixed-width
termination and weighted-degree bound.
For exact Hom between the degree-zero bundles,
use source length $L=7$ and target length $M=3$,
as in \ref{D:fixed-partial-resolution-unit-atlas}.
The canonical triangle has the same locally
free left end module as in that proof. Its
target error has cohomological degrees
$[-3,-2]$ and its source error $[4,8]$.
Thus degree-zero Hom of the separated-length
complexes is exactly ordinary Hom of $F,G$,
over every coefficient algebra. Compose
through the canonical inclusion of the
length-three target into its length-seven
complex. To test identity, precompose with
the actual vector bundle $P(g)\twoheadrightarrow F$
and compare with its canonical map. Its
two-column presentation $[T^h\to Q]$ can
be used in the same finite complexes.
This avoids a homotopy-identity condition
on an equal-length partial resolution.
Boundary variables give exact equality
and inverse witnesses. Composition is
lifted to the original Cech--Cartier model
and projected once, retaining its actual
homotopies.

At component normal order $j$, give transition
and map variables weight $\max(1,j)$ initially,
but give residual frame entries and their
inverses weight $m$. Positive-order coefficients
have additive degree in their normal order;
the bounded number of column factors add at most
a fixed multiple of $m$. For each fixed-length
arrow diagram every equation consequently has
degree $O(m)$ through depth $aa_i m$.
After rescaling the positive weights by $M/m$
and giving the residual variables weight $M$,
the bound is $KM$. One-step Frobenius multiplies
normal orders and coordinate degrees by the
fixed number $p$. Exactness of its complex
requires no coefficient-Frobenius flatness:
locally the frame diagonalizes $T^h\to Q$
into reductions modulo $z^b$, whose Frobenius
pullbacks are reductions modulo $z^{pb}$;
these have the same exact Cartier resolutions.
The locally split contractible pairs and the
locally free end modules likewise survive
every pullback. The fixed spectral translations
and root-character carries are retained in
these Frobenius diagrams.

It remains to glue the actual normal components,
not merely their reduced spectral curves.
At an original node $i,j$, write root coordinates
$x,y$, and put $k=aa_i m_i$, $l=aa_jm_j$.
Over every coefficient algebra $B$ there is
the exact fibre-product identity
\[
 B[x,y]/(x^ky^l)=B[x,y]/x^k
       \mathbin{\times}_{B[x,y]/(x^k,y^l)}B[x,y]/y^l.
                                                        \tag{D.TP.7}
\]
It holds by coefficient comparison of monomials,
also when $B$ has nilpotents. Localization and
etale pullback preserve it. Finite projective
module pasting along this fibre product is the
kernel construction of
\citep[Section15.6, Lemmas15.6.4--15.6.9,
especially Tag0D2J]{StacksTreePastingD}.
Its proof and the homogeneous kernel construction
give the equivariant assertion for the fixed
diagonalizable node group. There are no triple
intersections of distinct original components.

Choose homogeneous fibre frames on the rectangle
$B[x,y]/(x^k,y^l)$. The difference of two
scalar-weight-one characters is a character of
$\mu_{m_i}\times\mu_{m_j}$. Each matrix entry
therefore has exponents in one fixed pair of
residue classes modulo $m_i,m_j$. It has at most
\[
                       (aa_i)(aa_j)                  \tag{D.TP.8}
\]
coefficient slots. A gluing isomorphism and its
inverse use $O(r^2)$ variables and quadratic
matrix-product equations. Its products have
only the two bounded root carries. Evaluation
of a component presentation on the rectangle
uses its full perfect complex, including
(D.TP.6), so no false exactness after an
underived restriction of $Q$ is used.
Only the single marked support at this node
remains after restriction; all other $Q_x$
terms are zero. The resulting fixed-width
Hom complexes against a homogeneous free
rectangle bundle also have $O(r^2)$ invariant
slots by (D.TP.8). Their frame and inverse
witnesses therefore impose the actual
restriction identifications at this cost.
Assign these node variables weight $M$.
The fixed-width composition formulas already
proved give weighted degree $O(M)$ for their
compatibility with the component maps and
Frobenius arrows.

Milnor pasting now produces actual vector
bundles on $a\mathcal A$ with the specified
maps, over every coefficient algebra.
Conversely every Artinian family can be
framed on these finite local rectangles,
and its component presentations were covered
above. Proper-Quot openness of the residual
HN windows is the argument already proved
in \ref{D:weighted-fixed-determinant-normal-presentation};
the subbundle conditions are proper images
of determinantal failures. These are open
restrictions, not additional coefficient
variables. Summing over the fixed number
of original components and nodes proves
all the stated bounds and scope.
\end{proof}
